% =====================================================================
% paper06_assessment_separation_v77
% Parent: Paper06 v76 (preserved unchanged); local-only audit-repair derivative.
% v77 (2026-10-27): substantive claim, unit, verification, benchmark and reference repairs.
% No external release or deposit is authorized by this revision.
% Build: tectonic paper06_assessment_separation_v77.tex (elsarticle, authoryear)
% =====================================================================
\documentclass[preprint,11pt,authoryear]{elsarticle}

\usepackage{amsmath,amssymb,amsthm}
\usepackage{booktabs}
\usepackage{array}
\usepackage{tabularx}
\usepackage{graphicx}
\graphicspath{{./}{figs/}{paper06_assessment_separation_v77_src/figs/}}
\usepackage{enumitem}
\usepackage[colorlinks=true,urlcolor=blue,linkcolor=blue,citecolor=blue]{hyperref}
\emergencystretch=3em
\setlength{\tolerance}{2000}

\newtheorem{theorem}{Theorem}
\newtheorem{proposition}{Proposition}
\newtheorem{lemma}{Lemma}
\newtheorem{corollary}{Corollary}
\theoremstyle{definition}
\newtheorem{definition}{Definition}
\theoremstyle{remark}
\newtheorem{remark}{Remark}

\newcommand{\Etyp}{E_{\mathrm{typ}}}
\newcommand{\Eend}{E_{\mathrm{end}}}
\newcommand{\Ephys}{E_{\mathrm{tube,phys}}}
\newcommand{\Vtyp}{\mathcal{V}_{\mathrm{typ}}}
\newcommand{\Vweak}{\mathcal{V}_{\mathrm{weak}}}
\newcommand{\Vphys}{\mathcal{V}_{\mathrm{phys}}}
\newcommand{\Vend}{\mathcal{V}_{\mathrm{end}}}
\newcommand{\Wp}{W_+}
\newcommand{\FP}{\mathrm{FP}_{\mathrm{agg}}}
\newcommand{\Tube}{\mathrm{Tube}}
\newcommand{\Succ}{\mathrm{Succ}}
\newcommand{\End}{\mathrm{End}}
\newcommand{\conv}{\operatorname{conv}}
\newcommand{\Rp}{\mathbb{R}^n_+}
\newcommand{\Q}{\mathcal{Q}}
\newcommand{\act}[1]{\textsc{#1}}

\journal{Journal of Multi-Criteria Decision Analysis}

\begin{document}

\begin{frontmatter}

\title{Transition Certification with Composite Indices: The Acceptance Gap Between Compensatory and Noncompensatory Rules}

\author[1]{Amin Abaee\corref{cor1}}
\cortext[cor1]{Corresponding author. Email: \href{mailto:amin_abaee@ut.ac.ir}{amin\_abaee@ut.ac.ir}. ORCID: \href{https://orcid.org/0000-0002-0019-1842}{0000-0002-0019-1842}.}
\affiliation[1]{organization={Independent Researcher}, city={Tehran}, country={Iran}}

\begin{abstract}
Composite indices can certify transitions under weight-adaptive or common-plan rules. We formulate the comparison as a robust finite-horizon assessment with path-wise constraints and declared disturbances. The full-cone aggregate is lossless at a fixed trajectory, yet the plan/weight quantifiers need not commute. For finite menus in the stated margin-first class, the weight-adaptive accepted set is $\conv(\mathcal D)+\Rp$ and the common-plan set is the union of single-plan acceptance orthants; an abstract convexification of required margins closes that gap, while physical implementation requires a separate control assumption.

An exact rational witness has a nonempty gap interior: at its canonical point both normalized margins are $6/5$, so the equal-weight margin sum $s_1+s_2=12/5$ exceeds the witness's licensing threshold $2$; each weighting licenses at least one primitive plan, though no single primitive plan meets both typed floors. A closed-form reserve threshold separates rescuable from impossible states. Under the declared linear-flow blend, fractional control closes the action-indexed gap, while the specified shared-shock blend does not restore feasibility at the canonical point; the complete-excursion result is conditional on its union-of-tubes assumption. On the witness-defined diagonal slice, every positive-elasticity member of the stated CES family falsely certifies an exact interval; the Leontief limit is typed-safe under the declared collapse convention. The shared-shock gap survives the exact $1/2$-cut, while a distinct replace reading with identical plan- and weight-independent dips collapses the distinction (SI S4). Northern-cod values supply a historical unit comparison only; the benchmark is not a fitted fishery model.
\end{abstract}

\begin{keyword}
multi-criteria decision analysis \sep compensatory and noncompensatory aggregation \sep composite indicators \sep scalarization \sep robustness \sep sustainability assessment \sep viability
\end{keyword}

\end{frontmatter}

\section{Introduction}\label{sec:intro}

\subsection{Compensability as an acceptance question}

Whenever decision-makers assess a sustainability transition, they face a familiar question: can strength in one area make up for failure in another? Multi-attribute value theory builds trade-offs into its value structure \citep{KeeneyRaiffa1976}. Outranking methods, by contrast, limit compensation through concordance and veto conditions \citep{Roy1996,Vincke1992}; work on sustainable development has also emphasized the incommensurability of some criteria \citep{Munda2005,MartinezAlier1998}.

The same question appears in environmental assessment as the distinction between weak and strong sustainability. \emph{Weak sustainability} represents substitution among forms of capital through an aggregate that must not decline \citep{WorldBank2011,Neumayer2013,Boos2015}. \emph{Strong sustainability} treats selected stocks as separately binding \citep{Daly1990,Ekins2003}. Research on the indicator theory of these positions is well developed, as is the analysis of compensability in the multi-criteria methods used for sustainability assessment \citep{Hanley1999,Martinet2011,CairnsMartinet2014,DoyenGajardo2020,Usubiaga2025,Cinelli2014,Schar2025}.

Much of this work classifies aggregation at a fixed profile of criterion values—a snapshot. This paper asks a different question, about \emph{transitions}: when a plan unfolds over time under declared disturbances, and several criteria must hold along the entire path, when can a compensatory aggregate certify a transition that a noncompensatory assessment rejects? The question is not how weights should be chosen; methods for eliciting them are well developed \citep{Nardo2008,Becker2017}. Nor is it the familiar static limitation of weighted sums in recovering nonconvex parts of a Pareto front \citep{DasDennis1997}. The focus is \emph{acceptance semantics}: the rule that classifies a composite reading of a transition as accepted or rejected—the sorting problem of decision aiding \citep{Roy1996}.

\subsection{Two acceptance semantics}

To see where the two assessments can diverge, imagine the transition problem at state $z$. It tracks typed-margin coordinates $s_i$, each expressed after a declared positive reference scale has been chosen so that zero is its binding floor, and has a finite menu $\mathcal A$ of plans and a declared disturbance set $D$. The weights act on these normalized coordinates; they do not turn unlike physical quantities into a conservation law. Both assessments impose the same physical safe and destination sets $S^{\mathrm{phys}}$ and $G^{\mathrm{phys}}$. Exact-tube semantics checks the entire path, not just its endpoint: for every $d\in D$, a plan's full tube must remain in $S^{\mathrm{phys}}$, and its successor set must lie in $G^{\mathrm{phys}}$.

\begin{samepage}
Now consider the floor test. Write $s(p)=(s_1(p),\dots,s_n(p))$ for the typed-margin vector at state $p$. For each weight $w$ in the full nonnegative cone $\Wp=\Rp\setminus\{0\}$, let $E_w(z)$ be the menu plans that meet those shared requirements and also keep $w\cdot s(p)\ge0$ at every tube state and every successor state, for every $d\in D$. Let $\Etyp(z)$ be the menu plans that meet the same physical, destination, and disturbance requirements but keep each typed floor $s_i(p)\ge0$ at those same states. Thus the two assessments share physical feasibility, destination attainment, and disturbance coverage; only the test applied to the typed margins differs.
\end{samepage}

The key choice is simple: must one plan work for every weighting, or may each weighting have its own plan? Formally, the two acceptance rules differ only in the order of the quantifiers:
\begin{itemize}[leftmargin=2em]
\item \emph{Common-plan acceptance}: $\exists a\in\mathcal A\;\forall w\in\Wp:\ a\in E_w(z)$. One plan must work for every weighting.
\item \emph{Weight-adaptive acceptance}: $\forall w\in\Wp\;\exists a_w\in\mathcal A:\ a_w\in E_w(z)$. For each weighting, some plan works; it may be a different one.
\end{itemize}
The first rule implies the second, but the converse need not hold. Because $\Wp$ contains every coordinate weight, the full-cone identity is $\Etyp(z)=\bigcap_{w\in\Wp}E_w(z)$. Among plans meeting the shared requirements, a single plan passes every weight exactly when it satisfies each typed floor; equivalently, common-plan acceptance is $\exists a\in\Etyp(z)$. The gap, when it occurs, comes from the weight-adaptive rule's ability to use different plans for different weights. This quantifier gap can occur even for static alternative margin vectors; here the tubes and disturbances supply its transition-assessment setting. Table~\ref{tab:protocols} places these two rules among the four protocols considered in the paper.

\begin{table}[!htbp]
\centering\small
\caption{Four acceptance protocols for a composite index on a transition datum. Protocols 3 (common-plan robustness) and 4 (typed acceptance) have the same accepted-plan and accepted-state sets on the full cone (Proposition~\ref{prop:hierarchy}(ii)); they are shown separately to distinguish the common-plan and typed formulations. The separation is Protocol 2 against either equivalent form.}
\label{tab:protocols}
\begin{tabularx}{\textwidth}{>{\raggedright\arraybackslash}p{0.25\textwidth} >{\raggedright\arraybackslash}p{0.27\textwidth} X}
\toprule
Protocol & Quantifier form & Reading \\
\midrule
1. Fixed-weight assessment & $a\in E_w(z)$ at a declared $w$ & one index, one selected plan: ordinary fixed-weight certification \\
2. Weight-adaptive (per-weight) acceptance & $\forall w\ \exists a_w\in E_w(z)$ & for each weighting $w$, some plan $a_w$ is licensed; the plan may vary with $w$: the compensatory reading, $\Vweak=\bigcap_w\mathcal V_w$ \\
3. Common-plan robustness & $\exists a\ \forall w:\ a\in E_w(z)$ & one plan serves every weighting; equivalent to Protocol 4 on the full cone \\
4. Coordinate-wise (typed) acceptance & $\exists a:\ a\in\Etyp(z)$ & separately binding floors enforced path-wise; equivalent to Protocol 3 on the full cone \\
\bottomrule
\end{tabularx}
\end{table}

An analogous split appears in adjustable robust optimisation: a here-and-now decision is fixed before the relevant parameter is known, while a wait-and-see decision may depend on it. Here the assessment weight plays the role of that parameter \citep{BenTal2004}. This is an analogy for the order of the acceptance rule, not a claim about how real auditors or policymakers behave. The result concerns a certification protocol applied to a composite index, not the suitability of composite indices as descriptive summaries.

The witness makes the split concrete. \act{fast} risks the first floor, \act{slow} the second. \act{staged} protects both only when the reserve covers its bridge; \act{no-switch} misses the destination. At the action-indexed canonical point $(x,s_1,s_2)=(\tfrac12,\tfrac65,\tfrac65)$, $s_1+s_2=\tfrac{12}5\ge2$, so every $w\in\Wp$ licenses \act{fast} or \act{slow}, yet none of the four plans is typed-admissible. This verdict is specific to the stated menu and disturbance class.


\pagebreak
\subsection{Contributions}

The paper makes five contributions, stated in the order of the text.
\begin{enumerate}[leftmargin=2em,label=(\roman*)]
\item \emph{Localisation and hierarchy} (Section~\ref{sec:general}). On the full cone the typed operator is exactly the common-plan intersection, $\Etyp(z)=\bigcap_{w\in\Wp}E_w(z)$, and at any fixed trajectory the aggregate is lossless (Lemma~\ref{lem:cone}); the chain typed $\subseteq$ scalarized $\subseteq$ physical $\subseteq$ endpoint holds for every weight, and restricting a nonempty weight family can only enlarge the compensatory accepted set (Propositions~\ref{prop:hierarchy}--\ref{prop:monotone}). The quantifier separation itself is a plan-selection phenomenon and can occur for static alternatives; the finite-horizon paths specify how the margins are generated in the transition application.
\item \emph{A finite-menu geometric characterisation} (Theorem~\ref{thm:menu}). Under the explicit margin-first hypothesis $m_a(w)=w\cdot d_a$, weight-adaptive acceptance is $\conv(\mathcal D)+\Rp$ and common-plan acceptance is $\bigcup_a(d_a+\Rp)$. Convexifying required-margin vectors closes this abstract gap; physical convexification of controls is a separate implementability assumption. Lemma~\ref{lem:sync} identifies when the margin-first path requirement agrees exactly with scalarizing after the within-path maximum.
\item \emph{An explicit rational witness} (Theorem~\ref{thm:witness}). On a four-plan datum the gap is a region with nonempty interior, equal to the \emph{impossibility region} $I$; the aggregate-versus-floor discrepancy region splits at the reserve level $x=1$ into $I$ and a \emph{rescue set} $R$ that a reserve-financed staged plan already serves; exact thresholds $\rho_1=(2-s_1)/s_2$ and $\rho_2=s_1/(2-s_2)$ locate the licensing regions of the two principal plans, with coordinate-axis weights classified separately; and the minimal reserve increment that converts impossibility into rescue is $\kappa^*(z)=(1-x)\,\mathbf 1[z\notin\Vtyp]$ (Proposition~\ref{prop:rescue}). The separation persists under hold-prefix extension of the horizon; a separate two-stage construction shows one way a repair stage can erase it by enlarging the next-stage accepted sets while retaining the same final terminal rules (Proposition~\ref{prop:persist}). The quantifier gap is sufficient but not necessary for false acceptance: at the canonical point $(x,s_1,s_2)=(1/2,6/5,6/5)$ the single fixed weight $(1,1)$ already accepts a primitive plan, because $12/5\ge2$, while no primitive plan meets both typed floors (Remark~\ref{rem:fixedweight}).
\item \emph{What closes the gap and what does not} (Theorem~\ref{thm:blend}, Proposition~\ref{prop:timeshare}). Fractional blending of the primitive plans -- menu convexification at the level of control flows -- makes the typed accepted set of the action-indexed margin class coincide exactly with the compensatory region, so that gap is a property of the finite deterministic menu, not of either doctrine. The statement is class-specific: under the shared-shock class the blend window is empty at the canonical point ($12/5<5/2$), so blending does not close that gap (Proposition~\ref{prop:commonshock}). The complete-excursion result is conditional on a policy whose visited set is the union of both full primitive tubes; it does not characterize ordinary finite-rate switching. A distinct three-switch implementation keeps both canonical margins positive, so no general failure of finite switching is claimed.
\item \emph{Robustness of the separation to the aggregator and to the disturbance class} (Theorem~\ref{thm:ladder}, Proposition~\ref{prop:commonshock}). Replacing the linear index by a constant-elasticity-of-substitution (CES) aggregator with elasticity $\sigma$ changes the gap but does not remove it: the accepted rungs are nested, on the witness-defined diagonal gap slice, each positive-elasticity member falsely certifies an exact interval with critical floors $1<5/4<\sqrt2<\varphi<\sqrt3<\dots\to2$; therefore no positive elasticity is uniformly safe over the full gap, while the Leontief member ($\sigma=0$), which is the typed operator, is uniformly safe. Under a shared environmental shock striking every plan the gap survives, reduced by an exact $1/2$-cut, with the licensing thresholds unshifted on the surviving region; the reserve route then fails on an exact band of states.
\end{enumerate}
The analytic results are proved in Appendices~\ref{app:general}--\ref{app:witness} and SI~S2--S4. Finite checks support only the enumerated cases and do not establish continuum claims. The Lean formalization covers selected algebraic witness identities and CES rungs; the broader real-exponent, continuity, and topological arguments are supported by the written proofs, not claimed as machine-checked. Section~\ref{sec:benchmark} gives a synthetic Schaefer-style biomass illustration with an optional historical scale mapping, not an empirically calibrated fishery realization.

\paragraph{What is not claimed} No ranking of weak and strong sustainability as doctrines; no proof that any particular set of boundaries ought to be treated noncompensatorily; no empirical fit, forecast, validation, or performance result for a resource system (published DFO values are a limited scale/context anchor only); no infinite-horizon, stochastic or partially observed extension. The impossibility region is a negative certificate relative to the specified plan menu, and the separation is exhibited for a declared disturbance class.

\subsection{Position in the literature}

The backward constructions used here are typed instances of robust-predecessor, capture-basin and reach--avoid constructions \citep{Aubin1991,Aubin2011,SaintPierre1994,Lygeros1999,AubinCatte2002}; the hierarchy of Proposition~\ref{prop:hierarchy}(i) is constraint-set monotonicity of viability kernels \citep{Aubin1991,Frankowska1989}; viability methods are established in natural-resource management \citep{DeLaraDoyen2008,Doyen2012}. The plan/weight quantifier order is analogous to here-and-now versus wait-and-see decisions, but the assessment weight is not a realized disturbance and fractional physical control is not a randomization over plans. In composite-indicator practice the critique of weights as value judgements presented as measurement \citep{Hickel2020,Fischer2022} and the construction of indices without weights from Pareto fronts \citep{Gao2023} both concern the index; the result here concerns acceptance under a declared menu and rule. We do not assert priority over the elementary fact $\exists a\,\forall w\ne\forall w\,\exists a_w$, nor over the convex-hull geometry of multi-objective optimisation; Section~\ref{sec:related} states how those shared ingredients differ from the transition-certification question addressed here.
